\section{ベイズ推定による回帰分析}
\label{b26_bayesian_regression}

\subsection{授業内容}

\subsubsection{概要}
前回学んだベイズ統計学の基本原理をもとに，本コマではベイズ推定を回帰分析に適用する方法について学ぶ。回帰モデルの各パラメータ（切片と傾き）を確率変数として扱い，そのベイズ推定法と解釈について理解する。また，事後分布から得られる推定値の種類（EAP, MAP, MEDなど）と確信区間の意味，さらに複数のモデルを比較するための周辺尤度とベイズファクターの概念を学ぶ。これにより，頻度主義的な回帰分析とベイズ的アプローチの共通点と相違点を理解し，それぞれの利点を活かした統計的推論ができるようになることを目指す。

\subsubsection{コマ主題細目}
\begin{description}
	\item[回帰モデルのベイズ表現] 単回帰モデル$Y_i = \beta_0 + \beta_1 X_i + \varepsilon_i$を確率モデルとして表現し，ベイズ推定の枠組みで捉え直す。誤差項$\varepsilon_i$が正規分布$N(0, \sigma^2)$に従うと仮定する従来の回帰モデルでは，パラメータ$\beta_0$，$\beta_1$，$\sigma^2$は未知の固定値として扱われるが，ベイズアプローチではこれらのパラメータを確率変数として扱う。具体的には，各パラメータに事前分布を設定し（例えば$\beta_0$，$\beta_1$に正規分布，$\sigma^2$に逆ガンマ分布など），データが与えられた下での事後分布$p(\beta_0, \beta_1, \sigma^2 | data)$を求める。ベイズの定理により，この事後分布は事前分布と尤度関数の積に比例することを理解する。ベイズ回帰モデルは，パラメータの不確実性を直接モデル化できるというメリットを持つことを学ぶ。

	\item[事後分布からの推定値と確信区間] ベイズ推定では，事後分布から様々な推定値を導出できることを学ぶ。代表的な推定値として，事後平均（EAP: Expected A Posteriori），事後最頻値（MAP: Maximum A Posteriori），事後中央値（MED: Median）がある。EAPは事後分布の期待値であり，二乗損失関数の下で最適な推定値となる。MAPは事後分布のモードであり，最尤推定値と密接に関連する（無情報事前分布の場合，MAPと最尤推定値は一致する）。MEDは事後分布の中央値であり，絶対損失関数の下で最適となる。また，事後分布からは95\%確信区間（Credible Interval）を導出できるが，これは頻度主義的な95\%信頼区間とは解釈が異なることを理解する。確信区間は「パラメータが区間内に入る確率が95\%である」と直接解釈できる点が特徴である。さらに，最高密度区間（HDI: Highest Density Interval）という，事後分布の密度が最も高い領域に基づく区間についても学ぶ。

	\item[事前分布の選択と感度分析] 回帰モデルのベイズ推定における事前分布の役割と選択方法について学ぶ。無情報事前分布から強い情報を持つ事前分布まで，様々な選択肢とその影響を理解する。事前分布の選択は，サンプルサイズが小さい場合に特に重要であり，事前情報の強さによって事後分布が大きく変わる可能性がある。このため，異なる事前分布のもとでの推定結果の安定性を確認する「感度分析」の重要性を学ぶ。無情報事前分布を用いた場合でも，ベイズ推定は最尤推定とは異なる結果を与えることがあり（特に複雑なモデルや少数のデータの場合），その違いの原因と解釈について理解する。また，事前分布の選択と報告の透明性の重要性についても学ぶ。

	\item[モデル比較と周辺尤度] ベイズ統計学におけるモデル比較の方法として，周辺尤度（marginal likelihood）とベイズファクター（Bayes Factor）の概念を学ぶ。周辺尤度$p(data|Model)$は，あるモデルの下でデータが観測される確率であり，ベイズの定理における分母（正規化定数）に相当する。これは$p(data|Model) = \int p(data|\theta, Model)p(\theta|Model)d\theta$と表される。つまり，パラメータ空間全体にわたる尤度と事前分布の積の積分として計算される。二つのモデル$M_1$と$M_2$を比較する場合，ベイズファクター$BF_{12} = \frac{p(data|M_1)}{p(data|M_2)}$がモデル$M_1$の$M_2$に対する相対的な証拠の強さを表す。ベイズファクターは，帰無仮説検定における$p$値とは異なり，帰無仮説を積極的に支持する証拠を提供できる点が特徴である。複雑なモデルでは周辺尤度の計算が困難であるため，ネストしたモデルの比較にはサヴェージ・ディッキー比（Savage-Dickey ratio）などの近似法が用いられることを簡単に紹介する。
	      \begin{flushright}
		      $\to$ \textcite{Lee201709}のP.88--103，\textcite{Toyoda201506}のP.56--60.
	      \end{flushright}

\end{description}
\subsubsection{キーワード}
\begin{itemize}
	\item ベイズ回帰モデル
	\item 事後分布からの推定値（EAP, MAP, MED）
	\item 確信区間と最高密度区間（HDI）
	\item 周辺尤度とベイズファクター
	\item 事前分布の選択と感度分析
\end{itemize}
\subsection{授業情報}
\paragraph{コマの展開方法}
講義
\paragraph{標準シラバスにおける位置づけ}
科目番号5；心理学統計法；(1)心理学で用いられる統計手法；J より高度な記述や推定を目指して
\subsubsection{予習・復習課題}
\paragraph{予習}
前回学んだベイズ統計学の基本概念（ベイズの定理，事前分布，事後分布など）を復習しておくこと。また，第\ref{b06_Corr}，\ref{b07_regression}，\ref{b08_MultipleRegression}講で学んだ相関分析と回帰分析の基本についても再確認しておくとよい。特に，回帰モデルにおけるパラメータの意味，最小二乗法と最尤法による推定方法の違いについて理解を深めておくことで，ベイズ推定との比較がしやすくなる。

\paragraph{復習}
授業で学んだ内容を定着させるために，以下の点について理解を深めておくこと：

\begin{enumerate}
\item 回帰モデルをベイズ的な枠組みで表現し，事前分布と事後分布の関係を説明できるようにする
\item 事後分布から得られる様々な推定値（EAP, MAP, MEDなど）の違いと特徴を理解する
\item 頻度主義的な信頼区間とベイズ的な確信区間の解釈の違いを明確に説明できるようにする
\item 事前分布の選択が事後分布にどのように影響するかを理解し，適切な事前分布を選択する基準を考える
\item 周辺尤度とベイズファクターの計算方法と解釈を理解し，モデル比較に活用できるようにする
\end{enumerate}

また，簡単な例（例：少数のデータを用いた単回帰モデル）を自分で考え，異なる事前分布を設定した場合の事後分布の変化を考察してみるとよい。次回以降の授業ではJASPなどのソフトウェアを用いたベイズ推定の実践を行うため，今回の概念的理解をしっかりと固めておくことが重要である。